% Generated draft: requires mathematical and source review.
\documentclass[11pt]{article}
\usepackage[margin=25mm]{geometry}
\usepackage{fontspec}
\setmainfont{Latin Modern Roman}
\usepackage{amsmath,amssymb,mathtools,unicode-math}
\setmathfont{Latin Modern Math}
\usepackage{enumitem,needspace,xurl,hyperref}
\setlength{\parindent}{0pt}
\setlength{\parskip}{5pt}
\setlength{\emergencystretch}{3em}
\newcommand{\N}{\mathbb{N}}
\newcommand{\Z}{\mathbb{Z}}
\newcommand{\Q}{\mathbb{Q}}
\newcommand{\R}{\mathbb{R}}
\newcommand{\C}{\mathbb{C}}
\newcommand{\F}{\mathbb{F}}
\newcommand{\A}{\mathbb{A}}
\newcommand{\PP}{\mathbb{P}}
\newcommand{\E}{\mathbb{E}}
\newcommand{\eps}{\varepsilon}
\newcommand{\dd}{\,\mathrm{d}}
\begin{document}
\section*{UnsolvedMath ID: 20001083}
\textbf{Source problem code:} catalogue entry
\subsection{Definitions and mathematical statement}
T. Tao's source problem asks for a hypergraph regularity lemma for subsets of large relative density in pseudorandom sparse hypergraphs, motivated by a new proof of arbitrarily long arithmetic progressions in the primes and possible broader applications. The source gives a graph analogue: for $|A|=|B|=N$, a suitably ``sparsely $c(\epsilon,\delta)$-quasirandom'' host $G_0\subseteq A\times B$ and $G\subseteq G_0$ with $|G|>\delta|G_0|$, there are equitable partitions into $s,t<C(\epsilon,\delta)$ classes with at least $(1-\epsilon)st$ relatively $\epsilon$-regular pairs. The exact host and relative-regularity definitions are not supplied. The presenter suggests a correctly organized induction could be cleaner than existing hypergraph arguments; this is a methodological remark, not a theorem.

The attached aid reports that the arithmetic application is known under explicit linear-forms hypotheses, while the unrestricted host question remains underspecified. Its independent weak regularization can be stated on edge space as follows. Let $\mu$ be the normalized host probability measure, $0\le f\le1$ a relative edge-density function, $0<\epsilon<1$, and $\mathcal T$ a specified family of test sets. Call a positive-mass cell $C$ weakly regular if for every $T\in\mathcal T$ with $\mu(T\cap C)\ge\epsilon\mu(C)$,
\[
\left|\E_\mu(f\mid T\cap C)-\E_\mu(f\mid C)\right|\le\epsilon.
\]
A normalized conditional-energy increment gives a partition of edge space into at most $2^{\lceil\epsilon^{-4}\rceil}$ cells whose irregular cells have total host mass at most $\epsilon$. This explicitly states a compatible weak test-set convention; the summary omits its exact norm definition. Such cells need not arise from equitable vertex partitions or a coherent hypergraph-complex hierarchy, and this assertion alone is not a sparse counting lemma.

The aid's counting obstruction is the $3$-partite host on $\F_q^3$ with normalized weight
\[
\nu(x,y,z)=q\mathbf1_{x+y+z=0}.
\]
It has mean one and conditional expectation one on every proper coordinate subset. For a nontrivial additive character $\chi$, however, $\E\nu(x,y,z)\chi(x)\chi(y)\chi(z)=1$, while the dense product average is zero. For the nondegenerate two-edge cherry,
\[
\E_{x,y,z,z'}\nu(x,y,z)\nu(x,y,z')\mathbf1_{z\ne z'}=0,
\]
whereas replacing $\nu$ by the dense weight $1$ gives $1-q^{-1}$. Perfect proper-coordinate marginals therefore do not by themselves supply the stronger counting control needed for the original sparse-transference program. The example does not refute results with adequate linear-forms hypotheses.
\subsection{Short English statement}
Seek sparse relative hypergraph regularity with enough pseudorandomness for arithmetic applications. The supplied aid separates a weak host-mass partition theorem from counting, and gives a sparse host whose perfect lower-coordinate marginals still fail a two-edge count.
\subsection{Sources}
\begin{itemize}
\item Local UnsolvedMath registry, permanent ID $20001083$, source code catalogue entry; canonical \texttt{aim-combinatorics-notes.json}, \texttt{notes[207]}.
\item AIM workshop list \emph{Recent trends in additive combinatorics}, collected by Ernie Croot and Vsevolod F. Lev, problem 1.2 attributed to T. Tao: \url{https://aimath.org/WWN/additivecomb/additivecomb.pdf}. The extraction's repeated less-than-or-equal glyph is interpreted as epsilon in the regularity notation; a page header is not part of the problem. No PDF inspection is claimed.
\item Attached research attempt 1: partial result, run and registry status partially solved, rigor proved, novelty confidence low, classification PARTIAL-PROGRESS. It supplies no full reference for the known linear-forms application or weak-partition result. Its machine-generated, non-peer-reviewed claims require independent verification before citation.
\end{itemize}
\end{document}
